% Common project-advancement problem and value-gap foundation.
\section{Project Development and Advancement}
\label{sec:p05-project-advancement}

This section solves the ex ante advancement decision at a fixed conjectured
qualified manufacturing-service price \(p_m\). It uses the deterministic
optimized route value defined above. CMO market consistency and the feedback
that selects \(p_m^*\) are imposed in the capacity-market section below.

\subsection{Innovation-object boundary and planned-project technology}
\label{subsec:p05-technology}

Developer \(i\) chooses one common control \(x_i\geq0\). Its canonical meaning
is original-drug innovation investment / project-advancement intensity. It is
broader than pure clinical-development effort, but it is not patent
applications, patent-generating effort, basic research, or upstream scientific
discovery. There is no class-specific control \(x_{ig}\).

Predetermined project-advancement productivity \(a_i>0\) converts the control
into the expected measure of viable projects reaching route planning:
\begin{equation}
  \lambda_i^{\mathrm{plan}}=a_ix_i.
  \label{eq:p05-planned-intensity}
\end{equation}
The control has units \(\mathsf X\), and \(a_i\) has units
\(\mathsf P/\mathsf X\), so
\(\lambda_i^{\mathrm{plan}}\) has units \(\mathsf P\).

The advancement-cost function is
\begin{equation}
  C_X(x_i)
  =
  \frac{\kappa}{1+\nu}x_i^{1+\nu},
  \qquad
  \kappa>0,\quad \nu>0.
  \label{eq:p05-advancement-cost}
\end{equation}
Here \(\kappa\) has units
\(\mathsf C/\mathsf X^{1+\nu}\), so \(C_X\) is measured in currency
per developer and decision cohort. Its marginal cost is
\(C_X'(x_i)=\kappa x_i^\nu\). The boundary \(\nu=1\) gives a quadratic
advancement cost.

\subsection{Expected optimized project value}
\label{subsec:p05-expected-value}

For developer \(i\), institutional regime \(M\), and fixed conjectured price
\(p_m\), define
\begin{equation}
  \begin{aligned}
  \Omega_i(M,p_m)
  &=
  \mathbb E_F\!\left[
    W_i(q,m;M,p_m)
  \right]\\
  &=
  \int W_i(q,m;M,p_m)\,dF(q,m).
  \end{aligned}
  \label{eq:p05-expected-value}
\end{equation}
The subscript \(i\) includes the dependence of route value on developer
manufacturing capability \(k_i\). Because the
route-value maximum contains abandonment with value zero,
\(W_i(q,m;M,p_m)\geq0\). Under the stated integrability condition,
\[
  0\leq\Omega_i(M,p_m)<+\infty.
\]
This expectation replaces the old theoretical-baseline inclusive-value
object; no random-utility aggregation is used.

\subsection{Ex ante advancement problem}
\label{subsec:p05-problem}

The developer solves
\begin{equation}
  \max_{x_i\geq0}
  \left\{
    \beta a_ix_i\Omega_i(M,p_m)
    -
    \frac{\kappa}{1+\nu}x_i^{1+\nu}
  \right\}.
  \label{eq:p05-advancement-objective}
\end{equation}
The decision precedes the \((q,m)\) draw and downstream route choice.
Conditional on \(M,p_m,a_i,k_i\), the expected value
\(\Omega_i(M,p_m)\) is therefore fixed with respect to the individual
choice \(x_i\).

The necessary and sufficient Kuhn--Tucker conditions are
\begin{equation}
  \begin{gathered}
  x_i\geq0,\qquad
  \kappa x_i^\nu-\beta a_i\Omega_i(M,p_m)\geq0,\\
  x_i\left[
    \kappa x_i^\nu-\beta a_i\Omega_i(M,p_m)
  \right]=0.
  \end{gathered}
  \label{eq:p05-kkt}
\end{equation}
For an interior solution, the FOC is
\begin{equation}
  \kappa x_i^\nu
  =
  \beta a_i\Omega_i(M,p_m).
  \label{eq:p05-foc}
\end{equation}
For \(x_i>0\), the objective's second derivative is
\begin{equation}
  -\kappa\nu x_i^{\nu-1}<0.
  \label{eq:p05-soc}
\end{equation}
More generally, \(x_i^{1+\nu}\) is strictly convex on the nonnegative
domain for every \(\nu>0\). Hence the objective is strictly concave.
Its superlinear cost dominates the linear benefit as \(x_i\to\infty\), so a
unique maximizer exists.

The unique optimal common advancement intensity is
\begin{equation}
  x_i^*(M,p_m)
  =
  \left[
    \frac{\beta a_i}{\kappa}
    \Omega_i(M,p_m)
  \right]^{1/\nu}.
  \label{eq:p05-optimal-advancement}
\end{equation}
If \(\Omega_i=0\), the formula gives the unique corner \(x_i^*=0\).
If \(\Omega_i>0\), the solution is interior. Dimensionally,
\(\beta a_i\Omega_i/\kappa\) has units \(\mathsf X^\nu\), so the
solution has units \(\mathsf X\).

\subsection{Fixed-price binary MAH channel}
\label{subsec:p05-binary-channel}

At a fixed \(p_m\), deterministic route choice implies the pointwise weak ordering
\[
  W_i(q,m;1,p_m)
  \geq W_i(q,m;0,p_m).
\]
Taking expectations gives the finite policy comparison
\begin{equation}
  \Omega_i(1,p_m)-\Omega_i(0,p_m)
  =
  \int
  \left[
    W_i(q,m;1,p_m)
    -W_i(q,m;0,p_m)
  \right]dF(q,m)
  \geq0.
  \label{eq:p05-binary-channel}
\end{equation}
For positive \(\Omega_i\),
\[
  \frac{\partial x_i^*}{\partial\Omega_i}
  =
  \frac{1}{\nu}
  \left(\frac{\beta a_i}{\kappa}\right)^{1/\nu}
  \Omega_i^{1/\nu-1}
  >0.
\]
Thus, at fixed \(p_m\),
\[
  \Omega_i(1,p_m)>\Omega_i(0,p_m)
  \quad\Longleftrightarrow\quad
  x_i^*(1,p_m)>x_i^*(0,p_m).
\]
The strict inequality in expected value requires a positive-probability set
of projects in \(\mathcal C_i(p_m)\), where entrusted production strictly
improves the pre-reform maximum. Otherwise both
expected value and advancement can have a zero reform effect.

This is a binary finite comparison, not a derivative with respect to \(M\).
It is also a fixed-price result, not the equilibrium-price comparative static
through \(p_m^*\). The regime changes neither \(a_i\), \(\kappa\), \(\nu\),
the cost function, nor the project distribution. MAH affects
\(x_i^*\) only through \(\Omega_i\). Manufacturing capability \(k_i\) affects
\(\Omega_i\) through route values and sorting, not through
\(a_ix_i\) or \(C_X(x_i)\). No financing object enters the baseline
advancement FOC.

\subsection{Downstream project-value gap}
\label{subsec:p05-value-gap}

For a useful value decomposition, let \(B_i\) denote the continuation-value
benchmark before one additional viable project reaches route planning, and
let \(K_i\) denote value after it reaches that stage. Define
\begin{equation}
  K_i=B_i+\Omega_i,
  \qquad
  K_i-B_i=\Omega_i.
  \label{eq:p05-value-gap}
\end{equation}
Substituting this accounting identity into the interior FOC gives
\begin{equation}
  C_X'(x_i)
  =
  \beta a_i(K_i-B_i).
  \label{eq:p05-value-gap-foc}
\end{equation}
This is an Akcigit-style downstream project-value interpretation. The
auxiliaries \(B_i\) and \(K_i\) are not recursive state variables, and the
identity does not claim that MAH directly raises scientific-research
productivity.

\subsection{Scope boundary}
\label{subsec:p05-scope}

The baseline contains one common \(x_i\), one cost function \(C_X\), and one
expected optimized project value \(\Omega_i\). It contains no separate
research and development controls, no financing constraint on \(x_i\), no
borrowing choice, no finance market, no patent outcome, no
novelty-class-specific control, and no CMO market-clearing condition in this
block. Project advancement is ex ante; route choice and observed
holder--producer separation are downstream.
